% answers_short.tex —— 2008 年·普通高等学校招生考试·湖南卷·文科 · 简短答案速查
% 由 tools/make_short_answers.py 自 answers.tex 的 \gksolhead/\gkanswer 提取，
% 并与 2008/2008_数学真题_湖南文_答案要点.md 逐题比对后生成，勿手改（改源 MD 后重跑）。

\begin{gknotice}
\textbf{使用说明}：本节只给最终答案便于快速核对，完整推导见解析版。答案为编写者独立推导，\textbf{非官方标准答案}。
\end{gknotice}

\gkpart{选择题}
\begin{center}\begin{tabular}{ccccccccccc}
\toprule
题号 & 1 & 2 & 3 & 4 & 5 & 6 & 7 & 8 & 9 & 10 \\
\midrule
答案 & B & A & C & B & D & A & D & C & B & C \\
\bottomrule
\end{tabular}\end{center}

\gkpart{填空题}
\begin{itemize}[leftmargin=2.6em]
  \item[\textbf{11.}] $2$。
  \item[\textbf{12.}] $60$。
  \item[\textbf{13.}] $5$。
  \item[\textbf{14.}] $(x-1)^2+y^2=1$；$k=\pm\dfrac{\sqrt3}3$。
  \item[\textbf{15.}] $\dfrac{16}3$；$\left(\dfrac{28}3,28\right]$。
\end{itemize}

\gkpart{解答题}
\begin{itemize}[leftmargin=2.6em]
  \item[\textbf{16.}] （1）$\dfrac78$；（2）$\dfrac38$。
  \item[\textbf{17.}] （1）$T=2\pi$；（2）$\dfrac{2\sqrt6+2\sqrt2}5$。
  \item[\textbf{18.}] （1）证明见解析；（2）$60^\circ$。
  \item[\textbf{19.}] （1）$\dfrac{x^2}\lambda+\dfrac{y^2}{\lambda-4}=1$；（2）$\left(4,\dfrac{16}3\right]$。
  \item[\textbf{20.}] （1）$a_3=2$，$a_4=4$；$a_{2k-1}=k$，$a_{2k}=2^k$（$k\in\mathbf N^*$）；（2）不存在这样的 $k$。
  \item[\textbf{21.}] （1）证明见解析；（2）$(-5,-1)$。
\end{itemize}
